\section{Simulation and model details}
\label{app:model}

Table~\ref{tab:w2} gives the fixed mixing matrix of Equation~(\ref{eq:process}) (full row rank; the
weights of each observed variable have unit norm) and Table~\ref{tab:cae} the layers of the
convolutional autoencoder, both as implemented in the authoritative repository.

\begin{table}[!ht]
\centering
\caption[The fixed mixing matrix]{The fixed $4\times8$ mixing matrix $W_2$, shown transposed: row $k$
holds the weights of $\tanh(z_{t,k})$ in the four observed variables $x_{t,1},\dots,x_{t,4}$.}
\label{tab:w2}
\begin{tabular}{crrrr}
\toprule
$k$ & $x_{t,1}$ & $x_{t,2}$ & $x_{t,3}$ & $x_{t,4}$ \\
\midrule
1 & $-0.60447198$ & $-0.04842081$ & $-0.05518327$ & $-0.04100004$ \\
2 & $-0.06150649$ & $0.36085142$ & $0.45874929$ & $-0.84305316$ \\
3 & $-0.37179578$ & $-0.83557846$ & $-0.21939336$ & $-0.40588165$ \\
4 & $-0.01161746$ & $-0.11135410$ & $-0.17469105$ & $-0.11509976$ \\
5 & $0.05885824$ & $-0.18152197$ & $-0.55464952$ & $0.26627250$ \\
6 & $0.52361117$ & $0.30422580$ & $0.30578469$ & $0.07542546$ \\
7 & $0.45120458$ & $0.15832824$ & $-0.13652085$ & $0.16952036$ \\
8 & $-0.10604239$ & $0.07924137$ & $-0.53672635$ & $-0.06530930$ \\
\bottomrule
\end{tabular}
\end{table}

\begin{table}[!ht]
\centering
\caption[Layers of the convolutional autoencoder]{Layers of the convolutional autoencoder (kernel 3, stride 2 and padding 1 for every
convolution; output padding 1 for the transposed convolutions).}
\label{tab:cae}
\begin{tabular}{llrr}
\toprule
Stage & Layer & Output shape & Parameters \\
\midrule
Input & standardised window $X_t$ & $4\times32$ & -- \\
Encoder & Conv1d $4\to8$, ReLU & $8\times16$ & 104 \\
 & Conv1d $8\to16$, ReLU & $16\times8$ & 400 \\
 & flatten, Linear $128\to8$ (bottleneck $h_t$) & 8 & 1032 \\
Decoder & Linear $8\to128$, ReLU; unflatten & $16\times8$ & 1152 \\
 & ConvTranspose1d $16\to8$, ReLU & $8\times16$ & 392 \\
 & ConvTranspose1d $8\to4$ (reconstruction $\hat X_t$) & $4\times32$ & 100 \\
\midrule
Total & & & 3180 \\
\bottomrule
\end{tabular}
\end{table}

\clearpage
\section{Supplementary component diagnostics}
\label{app:diagnostics}

Figure~\ref{fig:recon} shows a typical in-control window and its frozen-CAE reconstruction: the
decoder reproduces the smooth level structure of each channel but not the high-frequency
variation, which therefore dominates $\SPE_t$. Figure~\ref{fig:response} compares the in-control
distributions of the standardised components with their distributions under three shifts. It is a
post-hoc descriptive diagnostic computed from the first 100 D5 trajectories at each shift with the
frozen models; it was not a pre-specified output of the experiment.

\begin{figure}[!ht]
\centering
\includegraphics[width=\linewidth]{diagnostic_reconstruction_example.pdf}
\caption[Example window and CAE reconstruction]{One in-control D3 window (trajectory 1, window
1001) and its reconstruction by the frozen CAE ($\SPE_t=96.7$, against a D3 mean of 86).}
\label{fig:recon}
\end{figure}

\begin{figure}[!ht]
\centering
\includegraphics[width=\linewidth]{diagnostic_component_response.pdf}
\caption[Component distributions in and out of control (post-hoc)]{Post-hoc descriptive
diagnostic: distributions of $s_t$ and $\ell_t$, standardised by the D3 moments, in control (D3)
and at $\delta=0.4$, 1.0 and 3.0 (first 100 D5 trajectories each). Mean shifts: $s_t$ 0.21, 0.75,
1.86 and $\ell_t$ 0.10, 0.64, 2.55 standard deviations.}
\label{fig:response}
\end{figure}

\clearpage
\section{Supplementary Phase II results}
\label{app:results}

Figure~\ref{fig:d4} is the repository's own summary of the D4 validation in
Table~\ref{tab:ic}. Tables~\ref{tab:d5}--\ref{tab:ablation} give the complete out-of-control
results; every entry is computed from the committed per-trajectory run lengths, and no D5 run
length was censored. Figure~\ref{fig:boxplots} shows run-length distributions of the primary pair
at four shifts.

\begin{figure}[!ht]
\centering
\includegraphics[width=0.7\linewidth]{results_d4_arl0_validation.pdf}
\caption[Independent in-control validation on D4]{D4 in-control ARL ($\pm1$ SE, 500 trajectories)
of the six charts against the target of 370 (repository figure).}
\label{fig:d4}
\end{figure}

\begin{table}[!ht]
\centering
\caption[Out-of-control ARL of all charts]{Restricted out-of-control ARL (SE) of all six charts across the shift grid (D5, 500
trajectories per shift).}
\label{tab:d5}
\small
\begin{tabular}{rrrrrrr}
\toprule
$\delta$ & $H(0.05)$ & $H(0.10)$ & $H(0.20)$ & $S(0.05)$ & $R(0.05)$ & $M(0.05)$ \\
\midrule
\tabinput{tables/d5_arl1_full.tex}
\bottomrule
\end{tabular}
\end{table}

\begin{table}[!ht]
\centering
\caption[Primary paired comparison in detail]{Primary comparison in detail: paired mean difference
$\bar d$ (SE), median run lengths, share of $M(0.05)$ trajectories signalling at the first
opportunity (no $H(0.05)$ trajectory did at any shift), mean delay in post-change observations
($\ARL_1+31$) and the ratio of the two charts' delays in both units.}
\label{tab:primary}
\small
\begin{tabular}{rrrrrrrrr}
\toprule
 & & \multicolumn{2}{c}{MRL} & & \multicolumn{2}{c}{Obs.\ delay} & \multicolumn{2}{c}{Ratio $H/M$} \\
\cmidrule(lr){3-4}\cmidrule(lr){6-7}\cmidrule(l){8-9}
$\delta$ & $\bar d$ (SE) & $H$ & $M$ & $P(\mathrm{RL}_M=1)$ & $H$ & $M$ & obs. & opp. \\
\midrule
\tabinput{tables/primary_contrast.tex}
\bottomrule
\end{tabular}
\end{table}

\begin{table}[!ht]
\centering
\caption[Paired comparison with the ablations]{Paired mean run-length differences (SE) of $H(0.05)$ against its ablations; negative
values mean that $H(0.05)$ signalled sooner.}
\label{tab:ablation}
\small
\begin{tabular}{rrr}
\toprule
$\delta$ & $\mathrm{RL}_H-\mathrm{RL}_S$ & $\mathrm{RL}_H-\mathrm{RL}_R$ \\
\midrule
\tabinput{tables/ablation_contrast.tex}
\bottomrule
\end{tabular}
\end{table}

\begin{figure}[!ht]
\centering
\includegraphics[width=0.8\linewidth]{results_primary_run_length_boxplots.pdf}
\caption[Run-length distributions of the primary pair]{Run-length distributions of $H(0.05)$ and
$M(0.05)$ at four shifts (log scale; repository figure).}
\label{fig:boxplots}
\end{figure}

\clearpage
\section{Reproducibility and provenance}
\label{app:repro}

The experiment is the committed state of the authoritative repository
\texttt{pipeline-thesis-independent} (branch \texttt{simplify-pipeline-code}, commit
\texttt{e62bdee}). Data generation uses master seed 795 with independent PCG64DXSM streams per
dataset, trajectory and shift; the CAE training seed is 3795; all numerical routes run in
double precision (network in single precision) on one CPU thread. The frozen preprocessing
statistics, CAE weights, OC-SVM state, D3 moments and control limits, and the per-trajectory
D4/D5 run lengths are committed artifacts, and the production run that produced them took 207 s.
Every figure and generated table in this report is rebuilt, without refitting or recalibration,
by \texttt{scripts/make\_all\_figures.py} in the report workspace, which reads these artifacts
read-only; the source and derivation of each figure are recorded in
\texttt{provenance/final\_figure\_provenance.md}. The post-hoc diagnostic audit cited in
Section~4.5 is a separate analysis of the same frozen experiment and is not part of its
pre-specified design.
